\chapter{Composite Pati--Salam potential: exact rank-one Hessian no-go}

Generalized inner fluctuation уже вывел composite fields из project finite
Dirac block. Теперь они не считаются независимыми fundamental variables:
spectral potential вычисляется после буквальной подстановки
\(Y(\phi,\Sigma_4)\) и \(H_R(\Delta)\) в physical half-trace конечного
оператора.

\section{Что известно и чего нет в литературе}

Первичная работа arXiv:1507.08161 перечисляет composite content
\begin{equation}
 \phi\sim(2_R,2_L,1),\qquad
 \Delta\sim(2_R,1_L,4),\qquad
 \Sigma_4\sim(1_R,1_L,15),
\end{equation}
но прямо оставляет полный анализ composite scalar potential будущей работе.
Поэтому ниже не импортируется готовый phenomenological ansatz: potential
вычисляется из уже проверенного project operator.

\section{High-scale reduction}

Чтобы сохранить weak \(SU(2)_L\) на Pati--Salam scale, необходимо
\begin{equation}
 \langle\phi\rangle=0.
\end{equation}
Но composite Yukawa block имеет вид
\begin{align}
 Y={}&(k^\nu\phi+k^e\widetilde\phi)\Sigma_4
 +(k^u\phi+k^d\widetilde\phi)(I_4-\Sigma_4).
\end{align}
Следовательно,
\begin{equation}
 \boxed{\phi=0\quad\Longrightarrow\quad Y=0}
\end{equation}
для любого \(\Sigma_4\). Weak-singlet adjoint на этом background является
точно flat и не может самостоятельно выбрать high-scale vacuum.

Остаётся right-Majorana block
\begin{equation}
 H_{\dot aI\dot bJ}
 =k^{\nu_R*}\Delta_{\dot aJ}\Delta_{\dot bI}.
\end{equation}
Константа \(k^{\nu_R}\) поглощается в normalization поля и не влияет на
Hessian signature.

\section{Exact physical half-traces}

Обозначим
\begin{equation}
 \rho=\Tr(\Delta^\dagger\Delta),
 \qquad
 \tau=\Tr[(\Delta^\dagger\Delta)^2].
\end{equation}
Прямое умножение полного 32-by-32 KO6 operator даёт
\begin{align}
 \frac12\Tr D_F^2&=\rho^2,\\
 \frac12\Tr D_F^4&=\tau^2.
 \label{eq:composite-ps-half-traces}
\end{align}
Factor \(1/2\) удаляет particle--antiparticle doubling. Сто random complex
\(2\times4\) matrices подтверждают identities
\eqref{eq:composite-ps-half-traces} с абсолютной ошибкой ниже
\(2\times10^{-10}\).

После положительного rescaling coefficients общий vacuum-equivalent
potential имеет нормированную форму
\begin{equation}
 \boxed{V_\Delta=-\rho^2+\tau^2.}
 \label{eq:composite-ps-potential}
\end{equation}

\section{Желаемый rank-one saddle}

Gauge transformations приводят \(\Delta\) к singular-value form. Standard
Model stabilizer требует rank one, поэтому берём
\begin{equation}
 \Delta_{\rm SM}
 =\begin{pmatrix}v&0&0&0\\0&0&0&0\end{pmatrix},
 \qquad
 v=2^{-1/4}.
\end{equation}
Это stationary point с
\begin{equation}
 V(\Delta_{\rm SM})=-\frac14.
\end{equation}
Exact Hessian на всех 16 real components \(\Delta\) имеет spectrum
\begin{equation}
 \boxed{
 8\sqrt2\ (1),\qquad
 0\ (9),\qquad
 -2\sqrt2\ (6).}
 \label{eq:composite-ps-rank-one-hessian}
\end{equation}
Девять zero modes совпадают с числом broken Pati--Salam generators, но шесть
оставшихся directions отрицательны. Поэтому это не vacuum, а saddle.

\section{Нежелательный stationary orbit}

Нижний stationary orbit имеет два равных singular values,
\begin{equation}
 \Delta_{(2)}
 =\begin{pmatrix}v&0&0&0\\0&v&0&0\end{pmatrix},
 \qquad
 V(\Delta_{(2)})=-1.
\end{equation}
Его Hessian spectrum равен
\begin{equation}
 16\sqrt2\ (1),\qquad
 8\sqrt2\ (3),\qquad
 0\ (12),
\end{equation}
то есть отрицательных directions нет. Однако rank-two condensate разрушает
four-color до subgroup с двумя color doublets и locking с \(SU(2)_R\), а не
до \(SU(3)_c\times U(1)_Y\). Энергетически предпочтённый orbit не является
Standard Model vacuum.

\section{Почему universal singlet снова не помогает}

Пусть connected или external singlet меняет только radial dependence:
\begin{equation}
 V=f(\rho)+b\tau^2,
 \qquad b>0.
\end{equation}
Для singular squares \((p,q)\) rank-one point имеет \(q=0\). Radial
stationarity требует
\begin{equation}
 f'(p)+4bp^3=0.
\end{equation}
Но coefficient при включении второго singular value равен
\begin{equation}
 \left.\frac{\partial V}{\partial q}\right|_{q=0}
 =f'(p)=-4bp^3<0.
\end{equation}
Следовательно, любая поправка, зависящая только от общей нормы
\(\Delta\), оставляет rank-growth instability.

\begin{theorem}[Composite rank-one no-go]
Минимальный constrained composite spectral Pati--Salam branch не имеет
устойчивого Pati--Salam-to-Standard-Model vacuum. Weak preservation делает
\(\Sigma_4\) flat, желаемый rank-one \(\Delta\)-orbit имеет шесть exact
negative Hessian modes, а global lower orbit имеет rank two и неправильный
unbroken subgroup. Norm-only singlet extension не меняет этот verdict.
\end{theorem}

\section{Конструктивное указание для улучшения}

Следующая модель должна породить не ещё один universal mass coefficient, а
rank-sensitive invariant. Минимальный target имеет вид
\begin{equation}
 c_{\det}\det(\Delta\Delta^\dagger)
\end{equation}
или эквивалентное coupling к independent connected adjoint, которое
различает rank one и rank two. Его coefficient должен быть выведен из
расширенного finite geometry; добавлять его вручную нельзя.

\section{Литература и воспроизводимость}

Composite field formulas и statement о незавершённом potential analysis
сверены с A.~H.~Chamseddine, A.~Connes, W.~D.~van Suijlekom,
\emph{Grand Unification in the Spectral Pati--Salam Model},
arXiv:1507.08161. Fundamental-branch vacuum obstruction сопоставлен с
H.~Karimi, \emph{Symmetry Breaking and Proton Decay in Spectral Pati--Salam
Model}, arXiv:1905.04533.

Exact symbolic Hessian и full-matrix trace audit сохранены в
\path{s2t_v4_pati_salam_composite_potential_hessian.py} и
\path{s2t_v4_pati_salam_composite_potential_hessian_results.json}.